% =====================================================================
%  1bat-ud1-7.tex  —  HORA 7: Gimnàstica d'inequacions (mecanització i verificació)
% =====================================================================
\horatitol{Sessió 7 — Gimnàstica d'inequacions: mecanització i verificació}%
{Guanyar destresa resolent inequacions de tot tipus, escrivint la solució en interval i comprovant-la.}

\subsection{Idea clau}

Ja coneixem les tres regles (i el parany del signe). Avui les automatitzem amb una
bateria d'inequacions de dificultat creixent. La rutina de cada exercici serà sempre
la mateixa: \textbf{resoldre}, \textbf{escriure la solució en forma d'interval} i
\textbf{comprovar-la} substituint un valor.

\begin{idea}
\textbf{Rutina de treball} per a cada inequació:
\begin{enumerate}[leftmargin=1.6em, itemsep=2pt]
  \item \textbf{Simplifica}: desfés parèntesis i, si hi ha fraccions, multiplica per
        treure els denominadors.
  \item \textbf{Agrupa}: les $x$ a un costat, els números a l'altre.
  \item \textbf{Aïlla} la $x$, \textbf{girant} el símbol si divideixes per un negatiu.
  \item \textbf{Escriu} la solució en interval i \textbf{comprova} amb un valor.
\end{enumerate}
\end{idea}

\subsection{Exemples}

\begin{exemple}[amb fraccions]
Resolem $\dfrac{x+1}{2}-1>\dfrac{x}{3}$. Multipliquem tot per $6$ (el mínim comú
denominador) per treure les fraccions:
\[
6\cdot\frac{x+1}{2}-6\cdot 1>6\cdot\frac{x}{3}
\;\Rightarrow\; 3(x+1)-6>2x.
\]
Desfem el parèntesi i agrupem:
\[
3x+3-6>2x \;\Rightarrow\; 3x-3>2x \;\Rightarrow\; 3x-2x>3 \;\Rightarrow\; x>3.
\]
Solució: $(3,+\infty)$.
\end{exemple}

\begin{exemple}[el terme en $x$ queda negatiu]
Resolem $4x-7\ge 7x+5$. Agrupem:
\[
4x-7x\ge 5+7 \;\Rightarrow\; -3x\ge 12.
\]
Dividim per $-3$ (negatiu) i \textbf{girem}:
\[
x\le \frac{12}{-3} \;\Rightarrow\; x\le -4.
\]
Solució: $(-\infty,-4]$.
\end{exemple}

\subsection{Exercicis resolts}

\begin{exresolt}[doble parèntesi amb signe negatiu]
Resol $-2(x-4)\le 3x-1$, dóna la solució en interval i comprova-la.
\solucio
Desfem el parèntesi (compte amb el signe del $-2$):
\[
-2x+8\le 3x-1.
\]
Agrupem les $x$ a la dreta i els números a l'esquerra (així el coeficient de $x$
queda positiu i evitem girar):
\[
8+1\le 3x+2x \;\Rightarrow\; 9\le 5x \;\Rightarrow\; \frac{9}{5}\le x.
\]
És a dir $x\ge\dfrac{9}{5}=1,8$. Interval $[1,8;\,+\infty)$. \textbf{Comprovació} amb
$x=2$: $-2(2-4)=4$ i $3\cdot 2-1=5$; $4\le 5$, cert.
\resposta{$x\ge 1,8$, interval $[1,8;\,+\infty)$.}
\end{exresolt}

\begin{exresolt}[verificar una solució proposada]
Per a la inequació $2x+3>11$, digues si $x=3$ és una solució vàlida i comprova-ho.
\solucio
Resolent: $2x>8 \Rightarrow x>4$. El valor $x=3$ \textbf{no} compleix $x>4$, així que
\emph{no} hauria de ser solució. Comprovem substituint a l'original: $2\cdot 3+3=9$,
i $9>11$ és \textbf{fals}. Confirmat: $x=3$ no és solució.
\resposta{No: $x=3$ dóna $9>11$, que és fals; la solució és $x>4$.}
\end{exresolt}

\subsection{Exercicis per fer}

\noindent\textit{Per a cadascuna: resol, escriu la solució en interval i comprova-la
amb un valor.}

\begin{exercicis}
\begin{exlist}
  \item $2x+3>11$
  \item $5-3x\le 2$
  \item $3(x-2)-5x\le 4x+6$
  \item $\dfrac{x+1}{2}-1>\dfrac{x}{3}$
  \item $-2(x-4)\le 3x-1$
  \item $4x-7\ge 7x+5$
  \item \textbf{(Validació)} Per a la inequació de l'exercici 3, digues si $x=0$ és
        solució i comprova-ho substituint.
\end{exlist}
\end{exercicis}

% ---------------------------------------------------------------------
\fullsolucions{Sessió 7}
\begin{solucions}
\begin{sollist}
  \item $2x>8 \Rightarrow x>4$, interval $(4,+\infty)$.
  \item $-3x\le -3$; girem: $x\ge 1$, interval $[1,+\infty)$.
  \item $3x-6-5x\le 4x+6 \Rightarrow -2x-6\le 4x+6 \Rightarrow -6x\le 12$; girem:
        $x\ge -2$, interval $[-2,+\infty)$.
  \item (Mult.\ per $6$) $3(x+1)-6>2x \Rightarrow 3x-3>2x \Rightarrow x>3$, interval
        $(3,+\infty)$.
  \item $-2x+8\le 3x-1 \Rightarrow 9\le 5x \Rightarrow x\ge 1,8$, interval
        $[1,8;\,+\infty)$.
  \item $-3x\ge 12$; girem: $x\le -4$, interval $(-\infty,-4]$.
  \item Solució de l'exercici 3: $x\ge -2$. El valor $x=0$ \textbf{sí} que la compleix
        ($0\ge -2$). Comprovació a l'original: $3(0-2)-5\cdot 0=-6$ i
        $4\cdot 0+6=6$; $-6\le 6$, cert.
\end{sollist}
\end{solucions}
